\vspace{0.5em}\noindent\textbf{Event Name}\par

AWS Study Tour: Enterprise Cloud Architectures and Industry Applications
featuring Cloud Kinetics and Renova Cloud

This program provided a practical learning opportunity at AWS and
included professional presentations from representatives of AWS, Cloud
Kinetics, and Renova Cloud.

\vspace{0.5em}\noindent\textbf{Time}\par

May 23, 2026.

\vspace{0.5em}\noindent\textbf{Location}\par

Amazon Web Services Vietnam Office, Bitexco Financial Tower, Ho Chi Minh
City.

The event was also livestreamed and recorded on the AWS Study Group
YouTube channel, allowing participants to attend remotely or review the
sessions later.

\vspace{0.5em}\noindent\textbf{Role}\par

Online participant.

I joined the event as a learner, followed the presentations delivered by
industry professionals, and studied how Cloud architectures are
designed, deployed, and operated in enterprise environments.

\vspace{0.5em}\noindent\textbf{Main Content}\par

The event introduced knowledge related to enterprise Cloud architectures
and demonstrated how AWS can be applied to practical business problems.

One important topic was the difference between a basic Cloud system and
a system that meets enterprise requirements. An enterprise architecture
must not only function correctly but also satisfy criteria such as:

\begin{itemize}
\tightlist
\item
  Scalability as the number of users or the volume of data increases.
\item
  High availability and reduced service interruption.
\item
  Data security and access control.
\item
  Backup and disaster-recovery capabilities.
\item
  Monitoring of performance, logs, and resource costs.
\item
  Compliance with internal requirements and industry-specific
  regulations.
\end{itemize}

The speakers also discussed the role of a Solution Architect in
translating business requirements into technical architectures. A
Solution Architect must understand customer needs, identify functional
and non-functional requirements, and select appropriate AWS services.

Another notable topic was the transition from a development environment
to production. A production system requires automated deployment,
monitoring and alerting, change-control mechanisms, and a recovery plan
in case a new release fails.

The presentations from Cloud Kinetics and Renova Cloud helped me better
understand how enterprises use Cloud technologies to modernize systems,
improve scalability, and support digital transformation. The event also
allowed students to learn from Solution Architects currently working in
the Cloud industry.

\vspace{0.5em}\noindent\textbf{Participation Evidence}\par

\href{https://www.youtube.com/live/FKtMkUqyny4?si=SdhD3d67wFy7qNr6}{Watch
the event recording on YouTube}

Personal evidence to be added:

\begin{itemize}
\tightlist
\item
  A screenshot taken while the event video was playing.
\item
  A screenshot of the YouTube viewing history.
\item
  A screenshot showing the participant's account information.
\item
  Personal notes created while watching the event.
\end{itemize}

\vspace{0.5em}\noindent\textbf{Lessons Learned}\par

The first lesson I learned was that Cloud architecture design should not
focus only on making an application run. Developers must also consider
security, performance, scalability, monitoring, disaster recovery, and
operational costs.

The second lesson was that an AWS service should only be used when it is
appropriate for the system requirements. Using more services does not
automatically result in a better architecture. A suitable architecture
should be simple, manageable, scalable, and aligned with actual needs.

The third lesson was that production environments must be prepared and
operated differently from development environments. Sensitive
information should not be stored directly in source code; data should be
backed up; and the system should provide logs, metrics, alarms, and
incident-response procedures.

\vspace{0.5em}\noindent\textbf{Personal Contribution}\par

The knowledge gained from the event helped me better understand the
roles of components such as API Gateway, backend servers, databases,
storage systems, and monitoring services in my internship project.

I summarized the enterprise architecture principles introduced during
the event and compared them with the architecture of the project
currently being developed. Through this comparison, I identified areas
that could be improved, including monitoring, backup, secret management,
scalability, and cost control.

I also applied the principles of system layering, restricted access
permissions, and resource monitoring when evaluating the project's
security and operational readiness.
